% concatenated from Bulletin_250821-Rocha.tex
\documentclass[leqno,12pt]{article} %leqno is the option to put formula numbers on the left side
\setlength{\textheight}{23cm}
\setlength{\textwidth}{17cm}
\setlength{\oddsidemargin}{0cm}
\setlength{\evensidemargin}{0cm}
\setlength{\topmargin}{0cm}
%
\usepackage{amsmath, amssymb}
\usepackage{amsthm} %theorem environment option
%
\renewcommand{\baselinestretch}{1.2} %changing the interline spacing
\renewcommand{\thefootnote}{} %footnote counter
\newcommand\supp{\mathop{\rm supp}}
\newcommand\esssup{\mathop{\rm ess \, sup}}
\newcommand\essinf{\mathop{\rm ess \, inf}}

%
%%%%%%%%% Theorem-like environments %%%%%%%%%%%
%
\theoremstyle{plain} %text of this environment is typesetted in italics
\newtheorem{theorem}{\indent\sc Theorem}[section]
\newtheorem{lemma}[theorem]{\indent\sc Lemma}
\newtheorem{corollary}[theorem]{\indent\sc Corollary}
\newtheorem{proposition}[theorem]{\indent\sc Proposition}
\newtheorem{claim}[theorem]{\indent\sc Claim}
%
\theoremstyle{definition} %text of this environment is typesetted in roman letters
\newtheorem{definition}[theorem]{\indent\sc Definition}
\newtheorem{remark}[theorem]{\indent\sc Remark}
\newtheorem{example}[theorem]{\indent\sc Example}
\newtheorem{notation}[theorem]{\indent\sc Notation}
\newtheorem{assertion}[theorem]{\indent\sc Assertion}
%
%If a theorem-like environment should not be numbered,
%add * after \newtheorem, and delete the counter option such as [theorem].
\newtheorem*{remark0}{\indent\sc Remark}
%
%%%%% Proof %%%%%
\renewcommand{\proofname}{\indent\sc Proof.}
%The following commands are available in the proof environment:
%\begin{proof}
%\end{proof}
%The end of a proof is marked with a square.
%%%%%%%%%%%%%%%%%%%%%%%%%%%%%%%%%%%%%%%%%
\makeatletter
%The following command is available for address and e-mail address:
%\address{<address>}{<E-mail address>}
%%%%%%%% definition of "\address" command %%%%%%%%%%%
\def\address#1#2{\begingroup
\noindent\parbox[t]{7.8cm}{%
\small{\scshape\ignorespaces#1}\par\vskip1ex
\noindent\small{\itshape E-mail address}%
\/: #2\par\vskip4ex}\hfill%
\endgroup}%
%%%%%%%%%%
\makeatother
%%%%%%%%%%%%%%%%%%%%%%%%%%%%%%%%%%%%%%%%%
%
\pagestyle{myheadings}
\markright{ } %title of the running head option
%
\title{The $H^p(\mathbb{Z}^n) - H^q(\mathbb{Z}^n)$ boundedness of \\ the discrete Riesz potential} %title of the paper
%
\author{
%
%\small{Dedicated to Professor Xxx Yyy on his sixtieth birthday} dedication if necessary
%
\textsc{Pablo Rocha} %names of authors
}
\date{} %leave empty
%
%%%%%%%%%%%%%%%%%%%%%%%%%%%%%%%%%%%%%%%%%

\begin{document}

\maketitle

%%%%%%%%%%%%%%% footnote %%%%%%%%%%%%%%%%
\footnote{ %2020 MSC numbers
2020 \textit{Mathematics Subject Classification}.
43A17, 42B30, 42B25.
}
\footnote{ %key words and phrases
\textit{Key words and phrases}:
Discrete Hardy Spaces; Atomic Decomposition; Discrete Riesz Potential
}
%\footnote{ %acknowledgment of support etc. if any
%$^{*}$Thanks.
%}
%%%%%%%%%%%%%%%%%%%%%%%%%%%%%%%%%%%%%%%%%



\begin{abstract}
In [J. Class. Anal., vol. 26 (1) (2025), 63-76], we proved that the discrete Riesz potential $I_{\alpha}$ is a bounded operator 
$H^p(\mathbb{Z}^n) \to H^q(\mathbb{Z}^n)$ for $\frac{n-1}{n} < p \leq 1$, $\frac{1}{q} = \frac{1}{p} - \frac{\alpha}{n}$ and 
$0 < \alpha < n$. In this note, we extend such boundedness on the full range $0 < p \leq 1$. 
\end{abstract}


\section{Introduction}

Let $0 < \alpha < n$, the Riesz potential $\mathcal{R}_{\alpha}$ on $\mathbb{R}^n$ is defined by
\begin{equation} \label{Ralfa}
(\mathcal{R}_{\alpha} f)(x) = \int_{\mathbb{R}^n} f(y) |x-y|^{\alpha-n} dy.
\end{equation}
The well known Hardy-Littlewood-Sobolev inequality establishes that if $1 < p < \frac{n}{\alpha}$ and 
$\frac{1}{q} = \frac{1}{p} - \frac{\alpha}{n}$, then the operator $\mathcal{R}_{\alpha}$ is bounded from $L^{p}(\mathbb{R}^n)$ into 
$L^{q}(\mathbb{R}^n)$. For $p=1$ and $q = n/(n-\alpha)$, $\mathcal{R}_{\alpha}$ is of weak-type $(1, q)$ (see \cite{Stein}). For the 
endpoint $p = \frac{n}{\alpha}$ is known that $\mathcal{R}_{\alpha}$ is not a bounded operator 
$L^{\frac{n}{\alpha}}(\mathbb{R}^n) \to L^{\infty}(\mathbb{R}^n)$, however it can see that 
$\mathcal{R}_{\alpha} : L^{\frac{n}{\alpha}}(\mathbb{R}^n) \to BMO(\mathbb{R}^n)$ is bounded. This last estimate can also be obtained 
by duality. Indeed, $\mathcal{R}_{\alpha} : H^1(\mathbb{R}^n) \to L^{\frac{n}{n-\alpha}}(\mathbb{R}^n)$ is bounded (see \cite{Weiss}), 
where $H^1(\mathbb{R}^n)$ is the $1$-Hardy space on $\mathbb{R}^n$. In \cite{Feff}, C. Fefferman and E. Stein proved that $(H^1)^{*} = BMO$, 
this remarkable result joint with the fact that $(L^{\frac{n}{n-\alpha}})^{*} = L^{\frac{n}{\alpha}}$ allow to obtain, by duality, 
the $L^{\frac{n}{\alpha}}(\mathbb{R}^n) \to BMO(\mathbb{R}^n)$ boundedness of $\mathcal{R}_{\alpha}$.

The real Hardy spaces $H^p(\mathbb{R}^n)$, $\frac{n-1}{n} < p \leq \infty$, were first introduced by E. Stein and G. Weiss in \cite{Weiss}. Therein, the authors describe the $H^p$ theory in terms of systems of conjugate harmonic functions, for it the theory was only developed on the range $\frac{n-1}{n} < p \leq 1$ (see \cite[\S 5.16]{St}). They also showed that $\mathcal{R}_{\alpha}$ is a bounded operator from 
$H^p(\mathbb{R}^n)$ into $H^q(\mathbb{R}^n)$, when $\frac{n-1}{n} < p < \frac{n}{\alpha}$ and $\frac{1}{q} = \frac{1}{p} - \frac{\alpha}{n}$. 
Afterward, C. Fefferman and E. Stein in \cite{Feff} introduced real variable methods to characterize the Hardy spaces by means of maximal functions on the full range $0 < p \leq \infty$. This second approach brought greater flexibility to the theory. For $1 < p \leq \infty$, 
one has that $H^{p}(\mathbb{R}^{n}) \cong L^{p}(\mathbb{R}^{n})$, $H^{1}(\mathbb{R}^{n}) \subset L^{1}(\mathbb{R}^{n})$ strictly, and for 
$0 < p < 1$ the spaces $H^{p}(\mathbb{R}^{n})$ and $L^{p}(\mathbb{R}^{n})$ are not comparable.

The spaces $H^{p}(\mathbb{R}^{n})$ can also be characterized by atomic decompositions. That is, every distribution $f \in H^{p}$ can be expressed by
\[
f = \sum_j \lambda_j a_j, 
\]
where the $a_j$'s are $p$ - atoms, $\{ \lambda_j \} \in \ell^{p}$ and $\| f \|^{p}_{H^{p}} \approx \sum_j | \lambda_j |^{p}$.
For $0 < p \leq 1$, an $p$ - atom is a function $a(\cdot)$ supported on a cube $Q$ such that
\[
\| a \|_{\infty} \leq |Q|^{-1/p} \,\,\  \text{and} \,\, \int x^{\alpha} a(x) dx = 0 \,\,\, \text{for all} \,\, |\alpha| \leq 
n \left( \frac{1}{p}-1 \right).
\]
Such decompositions were obtained by R. Coifman in \cite{Coif} for the case $n=1$ and by R. Latter in \cite{Latter} for the case $n \geq 1$. 
These decompositions, in principle, allow to study the behavior of certain operators on $H^{p}(\mathbb{R}^{n})$ by focusing one's 
attention on individual atoms (see \cite{Bownik}, \cite{Ricci}, \cite{Pablo}). For instance, using such atomic decomposition, it obtains that the Riesz potential $\mathcal{R}_{\alpha}$ is bounded from $H^p(\mathbb{R}^n)$ into $L^q(\mathbb{R}^n)$, for $0 < p \leq 1$ and 
$\frac{1}{q} = \frac{1}{p} - \frac{\alpha}{n}$. In \cite{Taible}, M. H. Taibleson and G. Weiss, by means of a molecular decomposition for members in $H^p(\mathbb{R}^n)$, proved the $H^p(\mathbb{R}^n) \to H^q(\mathbb{R}^n)$ boundedness for $\mathcal{R}_{\alpha}$ on the full range 
$0 < p \leq 1$ (see also \cite{Krantz}). This extends the result obtained by E. Stein and G. Weiss above mentioned. A study about the 
behavior of Riesz potential $\mathcal{R}_{\alpha}$ on others function spaces was made by E. Nakai in \cite{Nakai}.  
 
For $0 < \alpha < n$, the discrete counterpart of (\ref{Ralfa}) is the discrete Riesz potential $I_{\alpha}$ on $\mathbb{Z}^n$, which is 
defined by
\begin{equation} \label{Riesz potential}
(I_{\alpha}b)(j) = \sum_{i \in \mathbb{Z}^n \setminus \{ j \}} \frac{b(i)}{|i-j |^{n - \alpha}}, \,\,\,\,\,\, j \in \mathbb{Z}^n.
\end{equation}
As one would expect, (\ref{Riesz potential}) has a similar behavior to that of (\ref{Ralfa}). Indeed, in \cite[p. 288]{Hardy}, for $n=1$, 
$0 < \alpha < 1$, $1 < p < 1/\alpha$ and $\frac{1}{q} = \frac{1}{p} - \alpha$, G. H. Hardy et al established the $\ell^p(\mathbb{Z}) - \ell^q(\mathbb{Z})$ boundedness of $I_{\alpha}$ (for the case $n \geq 1$ see \cite[Proposition $(a)$]{Wainger}, other proof was given in \cite[Theorem 3.1]{PRocha}).

The theory for Hardy spaces on $\mathbb{Z}^n$ was developed by S. Boza and M. Carro in \cite{Carro} (see also \cite{Boza}). There, the authors
gave a variety of distinct approaches to characterize the discrete Hardy spaces $H^p(\mathbb{Z}^n)$ analogous to the ones given for the Hardy spaces $H^p(\mathbb{R}^n)$. They also described an atomic decomposition for elements in $H^p(\mathbb{Z}^n)$.

Y. Kanjin and M. Satake in \cite{Kanjin}, based on some results of \cite{Boza}, constructed a molecular decomposition for 
$H^p(\mathbb{Z})$ analogous to the ones given by M. Taibleson and G. Weiss in \cite{Taible} for the Hardy spaces $H^p(\mathbb{R}^n)$. With this framework, they obtain the Marcinkiewicz multiplier theorem on $H^p(\mathbb{Z})$ and proved the $H^p(\mathbb{Z}) \to H^q(\mathbb{Z})$ boundedness of discrete Riesz potential $I_{\alpha}$, for $n=1$, $0 < p < \alpha^{-1}$ and $\frac{1}{q} = \frac{1}{p} - \alpha$. The 
$H^{p}(\mathbb{Z}) \to \ell^{q}(\mathbb{Z})$ boundedness of (\ref{Riesz potential}), with $n=1$, was observed by the author in 
\cite[Remark 11]{Rocha}. 

In \cite{PRocha}, by means of the atomic characterization of $H^p(\mathbb{Z}^n)$ developed in \cite{Carro} and the boundedness of discrete fractional maximal operator, we proved that $I_{\alpha}$ is a bounded operator $H^{p}(\mathbb{Z}^n) \to \ell^{q}(\mathbb{Z}^n)$ for $n \geq 1$,
$0 < p < \frac{n}{\alpha}$ and $\frac{1}{q} = \frac{1}{p} - \frac{\alpha}{n}$. 

The $H^p(\mathbb{Z}^n) \to H^q(\mathbb{Z}^n)$ boundedness of $I_{\alpha}$ for $\frac{n-1}{n} < p \leq q \leq 1$, was proved by the author in \cite{Rocha2}. To achieve this result we used the characterization of $H^p(\mathbb{Z}^n)$ by the discrete Riesz transforms and furnished 
a molecular decomposition, as in \cite{Kanjin}, for the elements of $H^{p}(\mathbb{Z}^n)$ on the range $\frac{n-1}{n} < p \leq 1$ with 
$n \geq 1$. We point out that the lower bound $\frac{n-1}{n}$ in this decomposition is a consequence of Theorem 2.6 in \cite{Carro}.

The purpose of this note is to extend the $H^p(\mathbb{Z}^n) \to H^q(\mathbb{Z}^n)$ boundedness of $I_{\alpha}$ obtained in \cite{Rocha2} to the full range $0 < p \leq 1$.

The main result of this note is contained in the following theorem, which will be proved in Section 3 by means of the characterization 
maximal for $H^p(\mathbb{Z}^n)$ established in \cite[Theorem 2.7]{Carro} and the atomic decomposition for $H^p(\mathbb{Z}^n)$ also given in \cite{Carro}, joint with some auxiliary results of Section 2.

\

{\sc Theorem} \ref{main result}.
{\it For $0 < \alpha < n$, let $I_{\alpha}$ be the discrete Riesz potential given by (\ref{Riesz potential}). If
$0 < p \leq 1$ and $\frac{1}{q} = \frac{1}{p} - \frac{\alpha}{n}$, then
\[
\| I_{\alpha} \, b \|_{H^{q}(\mathbb{Z}^n)} \leq C \| b \|_{H^{p}(\mathbb{Z}^n)}
\]
for all $b \in H^{p}(\mathbb{Z}^n)$, where $C$ does not depend on $b$.}

\

{\bf Notation.} Throughout this paper, $C$ will denote a positive real constant not necessarily the same at each occurrence. We set 
$\mathbb{Z}_{+} = \mathbb{N} \cup \{0\}$. For every $A \subset \mathbb{Z}^n$, we denote by $\#A$ and $\chi_{A}$ the cardinality of 
the set $A$ and the characteristic sequence of $A$ on $\mathbb{Z}^n$ respectively. Given a real number $s \geq 0$, we write 
$\lfloor s \rfloor$ for the integer part of $s$. For $j=(j_1, ..., j_n) \in \mathbb{Z}^n$ we consider the two following norms 
$|j| = (j_1^2 + \cdot \cdot \cdot + j_n^2)^{1/2}$ and $|j|_{\infty} = \max \{ |j_k| : k = 1, ..., n \}$. Given $k^0 \in \mathbb{Z}^n$ and 
$m \in \mathbb{Z}_{+}$ fix, we put $Q_{k^0, m} = \{ i \in \mathbb{Z}^n : |i - k^0|_{\infty} \leq m \}$, that is the discrete cube centered 
at $k^0$ and side length $2m + 1$. If $\beta$ is the multiindex $\beta =(\beta_{1},...,\beta _{n}) \in \mathbb{Z}_{+}^{n}$, then 
$[\beta] :=\beta _{1}+...+\beta _{n}$ and $j^{\beta} := j_{1}^{\beta_1} \cdot \cdot \cdot j_{n}^{\beta_n}$ for every $j=(j_1, ..., j_n) \in \mathbb{Z}^n$. By $B_r({\bf 0 })$ we denote the Euclidean ball in $\mathbb{R}^n$ with center at ${\bf 0 }$ and radius $r > 0$. Given a function $F$ defined on $\mathbb{R}^n$, we shall use the notation $F^d$ to indicate its restriction on $\mathbb{Z}^n$. 

\section{Preliminaries}

The Schwartz space $\mathcal{S}(\mathbb{R}^{n})$ is defined by
\[
\mathcal{S}(\mathbb{R}^{n}) = \left\{ \phi \in C^{\infty}(\mathbb{R}^{n}) : \sup_{x \in \mathbb{R}^{n}} (1+|x|)^{N} 
|(\partial^{\beta}\phi)(x)| < \infty \,\,\, \forall \,\, N \in \mathbb{Z}_{+}, \, \beta \in \mathbb{Z}_{+}^{n}  \right\}.
\]
We topologize the space $\mathcal{S}(\mathbb{R}^{n})$ with the following family of seminorms
\[
\| \phi \|_{\mathcal{S}(\mathbb{R}^{n}), \, N} = \sum_{[\beta] \leq N} \sup_{x \in \mathbb{R}^{n}} (1+|x|)^{N} |(\partial^{\beta} \phi)(x)|, \,\,\,\,\,\,\, (N \in \mathbb{Z}_{+}).
\]

For $0 < p < \infty$ and a sequence $b = \{ b(i) \}_{i \in \mathbb{Z}^n}$ we say that $b$ belongs to $\ell^{p}(\mathbb{Z}^n)$ if
\[
\| b \|_{\ell^p(\mathbb{Z}^n)} :=\left( \sum_{i \in \mathbb{Z}^n} |b(i)|^p \right)^{1/p} < \infty.
\]
For $p=\infty$, we say that $b$ belongs to $\ell^{\infty}(\mathbb{Z}^n)$ if 
\[
\|b \|_{\ell^\infty(\mathbb{Z}^n)} := \sup_{i \in \mathbb{Z}^n} |b(i)| < \infty.
\]

Given two sequences $b=\{ b(i) \}_{i \in \mathbb{Z}^n}$ and $c=\{ c(i) \}_{i \in \mathbb{Z}^n}$, their convolution $b \ast c$ is defined by
\[
(b \ast c)(j) = \sum_{i \in \mathbb{Z}^n} b(i) \, c(j-i),
\]
when the series converges (e.g. \cite[Section 1.2.3, p. 21-25]{Loukas}). Moreover, $b \ast c = c \ast b$.

We recall that a discrete cube $Q$ centered at $j =(j_1, ..., j_n)  \in \mathbb{Z}^n$ can be written of the form 
$Q = Q_{j, m} = \prod_{1 \leq l \leq n} [[j_l -m, j_l +m]]$, where for each $l=1, ..., n$, $[[j_l -m, j_l +m]] := [j_l -m, j_l +m] \cap \mathbb{Z} = \{ j_l - m, ..., j_l, ..., j_l + m \}$ with $m \in \mathbb{Z}_{+}$. It is clear that $\# Q = (2m+1)^n$.

Let $0 \leq \alpha < n$, given a sequence $b = \{ b(i) \}_{i \in \mathbb{Z}^n}$ we define the centered fractional maximal sequence 
$M_{\alpha} b$ by
\[
(M_{\alpha} b)(j) = \sup_{Q \ni j} \frac{1}{(\# Q)^{1 - \frac{\alpha}{n}}} \sum_{i \in Q} |b(i)|, \,\,\,\,\,\, j \in \mathbb{Z}^n,
\]
where the supremum is taken over all discrete cubes $Q$ centered at $j$. We observe that if $\alpha = 0$, then $M_0 = M$ where $M$ is the centered discrete maximal operator of Hardy-Littlewood. For $0 \leq \alpha < n$, by Theorem 2.3 and Proposition 2.4 in \cite{PRocha}, we have that the operator $M_{\alpha}$ is bounded from $\ell^{p}(\mathbb{Z}^n)$ into $\ell^{q}(\mathbb{Z}^n)$ for $1 < p < \frac{n}{\alpha}$ and 
$\frac{1}{q} = \frac{1}{p} - \frac{\alpha}{n}$.

We adopt the following definition of discrete Hardy space on $\mathbb{Z}^n$ (see \cite[Theorem 2.7]{Carro}). Let $\Phi \in \mathcal{S}(\mathbb{R}^n)$ with $\int_{\mathbb{R}^n} \Phi =1$. Then, for $t>0$, we put $\Phi_t^d(j) = t^{-n} \Phi(j/t)$ if $j \in \mathbb{Z}^n \setminus \{ {\bf 0} \}$ and $\Phi_t^d({\bf 0}) = 0$. Now, for $0 < p \leq \infty$, we define
\begin{equation} \label{discrete Hp}
H^p(\mathbb{Z}^n) := \left\{ b \in \ell^p(\mathbb{Z}^n) : \sup_{t>0} |(\Phi_t^d \ast b)| \in \ell^p(\mathbb{Z}^n) \right\},
\end{equation}
equipped with the quasi-norm given by
\begin{equation} \label{discrete norm}
\| b \|_{H^p(\mathbb{Z}^n)} := \| b \|_{\ell^p(\mathbb{Z}^n)} + \|\sup_{t>0} |(\Phi_t^d \ast b)| \|_{\ell^p(\mathbb{Z}^n)}.
\end{equation}

In \cite{Carro}, Boza and Carro also gave an atomic characterization of $H^p(\mathbb{Z}^n)$ for $0 < p \leq 1$. Before establishing this result we recall the definition of $(p, \infty, N_p)$-atom in $H^p(\mathbb{Z}^n)$. 

\begin{definition} Let $0 < p \leq 1$ and $N_p := \lfloor n(p^{-1} - 1) \rfloor$. We say that a sequence 
$a = \{ a(j) \}_{j \in \mathbb{Z}^n}$ is an $(p, \infty, N_p)$-atom centered at a discrete cube $Q \subset \mathbb{Z}^n$ if the following three conditions hold:

(a1) $\supp a \subset Q$,

(a2.$p$) $\| a \|_{\ell^\infty(\mathbb{Z}^n)} \leq (\# Q)^{-1/p}$,

(a3.$p$) $\displaystyle{\sum_{j \in Q}} j^{\beta} a(j) = 0$ for every multi-index $\beta=(\beta_1, ..., \beta_n) \in \mathbb{Z}_{+}^{n}$ with 
$[\beta] \leq N_p$.
\end{definition}

The atomic decomposition mentioned for $H^p(\mathbb{Z}^n)$ is established in the following theorem.

\begin{theorem} (\cite[Theorem 3.7]{Carro}) \label{atomic Hp} Let $0 < p \leq 1$, $N_p = \lfloor n (p^{-1} - 1) \rfloor$ and 
$b \in H^{p}(\mathbb{Z}^n)$. Then there exist a sequence of $(p, \infty, N_p)$-atoms $\{ a_k \}_{k=0}^{+\infty}$, a sequence of scalars 
$\{ \lambda_k \}_{k=0}^{+\infty}$ and a positive constant $C$, which depends only on $p$ and $n$, with 
$\sum_{k=0}^{+\infty} |\lambda_k |^{p} \leq C \| b \|_{H^{p}(\mathbb{Z}^n)}^{p}$ such that $b = \sum_{k=0}^{+\infty} \lambda_k a_k$, where the series converges in $H^{p}(\mathbb{Z}^n)$.
\end{theorem}

\begin{remark} \label{Hp equiv}
From \cite[Corollary 4]{Rocha}, it follows for every $0 < p \leq 1$ that $H^p(\mathbb{Z}^n) \subset \ell^p(\mathbb{Z}^n)$ strictly. If 
$\Phi \in \mathcal{S}(\mathbb{R}^n)$ is such that $\supp(\Phi) \subset B_{1/4}({\bf 0 })$, then 
$\sup_{t>0} |(\Phi_t^d \ast b)(j)| \leq C (Mb)(j)$ for all $j \in \mathbb{Z}^n$, where $M$ is the discrete maximal operator of 
Hardy-Littlewood. Thus, by \cite[Theorem 2.3]{PRocha}, $H^p(\mathbb{Z}^n) = \ell^p(\mathbb{Z}^n)$ for $1 < p \leq \infty$, with equivalent norms. 
\end{remark}

We conclude these preliminaries with the following two auxiliary results.

\begin{lemma} For $0 < \alpha < n$, $m \in \mathbb{N}$ and $k^0, k \in \mathbb{Z}^n$ fix, the estimate
\begin{equation} \label{kernel estimate}
\sum_{i \in Q_{k^0, m} \setminus \{ k \}} |k-i|^{\alpha - n} \leq C (2m+1)^{\alpha}
\end{equation}
holds with $C$ independent of $m \in \mathbb{N}$ and $k^0, k \in \mathbb{Z}^n$.
\end{lemma}

\begin{proof}
Since $|j|_{\infty} \leq |j|$ for all $j \in \mathbb{Z}^n$, we have
\[
\sum_{i \in Q_{k^0, m} \setminus \{ k \}} |k-i|^{\alpha - n} \leq \sum_{i \in Q_{k^0, m} \setminus \{ k \}} |k-i|^{\alpha - n}_{\infty}
= \sum_{i \in I_{k^0, k, m}} |k-i|^{\alpha - n}_{\infty} + \sum_{i \in J_{k^0, k, m}} |k-i|^{\alpha - n}_{\infty},
\]
where
\[
I_{k^0, k, m} = \{ i \in Q_{k^0, m} : |k-i|_{\infty} > m \} \,\,\, \text{and} \,\,\, J_{k^0, k, m} = \{ i \in Q_{k^0, m} : 0<|k-i|_{\infty} 
\leq m \}.
\]
Now, one has
\[
\sum_{i \in I_{k^0, k, m}} |k-i|^{\alpha - n}_{\infty} \leq m^{\alpha - n} \cdot \#Q_{k^0, m} \leq C (2m+1)^{\alpha}.
\]
On the other hand, there exists a unique $s_0 \in \mathbb{N}$ such that $2^{s_0 - 1} \leq m < 2^{s_0}$. Then,
\[
\sum_{i \in J_{k^0, k, m}} |k-i|^{\alpha - n}_{\infty} \leq \sum_{s=0}^{s_0 - 1} \sum_{2^{-(s+1)} m < |k-i|_{\infty} \leq 2^{-s} m} 
|k-i|^{\alpha - n}_{\infty}
\]
\[
\leq m^{\alpha} \sum_{s=0}^{s_0 - 1} 2^{-(s+1) \alpha} \, \left( (2^{-(s+1)} m)^{-n} \sum_{i: |k-i|_{\infty} \leq 2^{-s} m} 1 \right)
\]
\[
\leq C (2m+1)^{\alpha} \sum_{s=0}^{\infty} 2^{-(s+1) \alpha}.
\]
Thus, (\ref{kernel estimate}) follows.
\end{proof}

\begin{lemma} \label{phimax}
Let $0 < \alpha < n$ and $\Phi \in \mathcal{S}(\mathbb{R}^n)$ such that $\supp(\Phi) \subset B_{1/4}({\bf 0})$. If $I_{\alpha}$ is the discrete Riesz potential given by (\ref{Riesz potential}), then there exists a positive constant $C = C_{\Phi, n, \alpha}$ such that
\begin{equation} \label{ineq puntual}
\sup_{t>0} |(\Phi_t^d \ast I_{\alpha}a)(j)| \leq C (\# Q)^{\frac{\alpha}{n} - \frac{1}{p}},
\end{equation}
holds for all $j \in \mathbb{Z}^n$ and all $p$-discrete atoms $a(\cdot)$ supported on $Q$.
\end{lemma}

\begin{proof}
Let $a(\cdot)$ be a $p$-atom supported by the cube $Q_{k^0, m} = \{ i : | i - k^0 |_{\infty} \leq m \}$ centered at $k^0$ and 
side length $2m+1$. Since by hypothesis $\supp(\Phi) \subset B_{1/4}({\bf 0})$ and $\Phi^{d}_{t}({\bf 0}) = 0$, we have
\[
(\Phi_t^d \ast I_{\alpha}a)(j) = 0, 
\]
for all $0 < t < 1$ and $j \in \mathbb{Z}^n$.

For $t \geq 1$ we have $1 \leq \lfloor t \rfloor \leq t$, then the condition (a2.$p$) of the atom $a(\cdot)$ and the estimate 
(\ref{kernel estimate}) give
\begin{align*}
|(\Phi_t^d \ast I_{\alpha}a)(j)| &\leq \sum_{k} |\Phi_{t}(j-k)| \sum_{i \in Q_{k^0, m} \setminus \{ k \}} |a(i)| |k-i|^{\alpha - n} \\
&\leq C (\# Q)^{\frac{\alpha}{n} - \frac{1}{p}} \left( t^{-n} \sum_{|k|_{\infty} \leq \frac{t}{4}} |\Phi(k/t)| \right) \\
&\leq C (\# Q)^{\frac{\alpha}{n} - \frac{1}{p}} \left( \lfloor t \rfloor^{-n} \sum_{|k|_{\infty} \leq \lfloor t \rfloor} 
|\Phi(k/t)| \right) \\
&\leq C \|\Phi \|_{\infty} (\# Q)^{\frac{\alpha}{n} - \frac{1}{p}},
\end{align*}
for all $1 \leq t < \infty$ and $j \in \mathbb{Z}^n$. Thus, (\ref{ineq puntual}) follows.
\end{proof}

\section{Main result}

In this section we establish the $H^p(\mathbb{Z}^n) - H^q(\mathbb{Z}^n)$ boundedness of the discrete Riesz potential $I_{\alpha}$ on the full range $0 < p \leq 1$, where $\frac{1}{q} = \frac{1}{p} - \frac{\alpha}{n}$. For them, we first show that the operator $I_{\alpha}$ is bounded uniformly in the $H^q(\mathbb{Z}^n)$-norm on all $(p, \infty, N_p)$-atoms.

\begin{proposition} \label{maximal atom estim}
Let $0 < \alpha < n$ and $\Phi \in \mathcal{S}(\mathbb{R}^n)$ such that $\supp(\Phi) \subset B_{1/4}({\bf 0})$. If $I_{\alpha}$ is the discrete Riesz potential given by (\ref{Riesz potential}), then, for $0 < p \leq 1$ and $\frac{1}{q} = \frac{1}{p} - \frac{\alpha}{n}$, there exists a positive constant $C$ such that
\begin{equation} \label{uniform estimate}
\|\sup_{t>0}|(\Phi_{t}^d \ast I_{\alpha} a)| \|_{\ell^{q}(\mathbb{Z}^n)} \leq C,
\end{equation}
for all discrete $(p, \infty, N_p)$-atom $a = \{ a(i) \}$.
\end{proposition}

\begin{proof}
To prove (\ref{uniform estimate}), let $a(\cdot)$ be a discrete $(p, \infty, N_p)$-atom centered at the cube 
$Q_{k^0}= \displaystyle{\prod_{1 \leq l \leq n}}\left[\left[ k^0_l - m, k^0_l + m \right]\right]$. We put
$4\lfloor \sqrt{n} \rfloor Q_{k^0} = \displaystyle{\prod_{1 \leq l \leq n}} \left[\left[ k^0_l - 4\lfloor \sqrt{n} \rfloor m, k^0_l + 4\lfloor \sqrt{n} \rfloor m \right]\right]$. Then, we split
\begin{align}
\sum_{j \in \mathbb{Z}^n} |\sup_{t>0}|(\Phi_{t}^d \ast I_{\alpha} \, a)(j)|^{q} &= \sum_{j \in 4\lfloor \sqrt{n} \rfloor Q_{k^0}} 
|\sup_{t>0}|(\Phi_{t}^d \ast I_{\alpha} \, a)(j)|^{q} \label{sum2} \\
&+ \sum_{j \in \mathbb{Z}^n \setminus 4\lfloor \sqrt{n} \rfloor Q_{k^0}} |\sup_{t>0}|(\Phi_{t}^d \ast I_{\alpha} \, a)(j)|^{q}. \nonumber
\end{align} 
To estimate the first sum in (\ref{sum2}), we apply Lemma \ref{phimax} above, which leads to
\begin{align}
\sum_{j \in 4\lfloor \sqrt{n} \rfloor Q_{k^0}}|\sup_{t>0}|(\Phi_{t}^d \ast I_{\alpha} \, a)(j)|^{q} &\leq  
C (\# Q_{k^0}) (\# Q_{k^0})^{q(\frac{\alpha}{n} - \frac{1}{p})} = C, \label{estim C}
\end{align}
with $C$ independent of $k^0$ and $m$.

To estimate the second sum in (\ref{sum2}), we analize the cases $0 < t < 1$ and $t \geq 1$ separately. For $0 < t < 1$, by taking into account that $\supp(\Phi) \subset B_{1/4}({\bf 0})$ and $\Phi^{d}_{t}({\bf 0}) = 0$, we have 
\begin{equation} \label{cero}
(\Phi_t^d \ast I_{\alpha}a)(j) = 0
\end{equation}
for all $j \in \mathbb{Z}^n$. For $t \geq 1$ we consider two sub-cases,

\

{\bf Case I:} $|j - k^0 |_{\infty} > 4 \lfloor \sqrt{n} \rfloor m$ and $1 \leq t \leq 2 |j - k^0 |_{\infty}$. We have
\[
(\Phi_t^d \ast I_{\alpha}a)(j) = \sum_{k : |j-k|_{\infty} \leq \frac{t}{4}} \Phi^{d}_{t}(j-k) \sum_{i \in Q_{k^0}} a(i) |i-k|^{\alpha-n}.
\]
Now, $|j-k|_{\infty} \leq t/4 \leq |j-k^0|_{\infty}/2$ and $|j - k^0 |_{\infty} > 4 \lfloor \sqrt{n} \rfloor m$ imply that 
$|k - k^0 |_{\infty} \geq |j - k^0 |_{\infty}/2 \geq 2|i - k^0 |_{\infty}$. We put $N - 1 = \lfloor n(p^{-1} - 1) \rfloor$. In view of the moment condition (a3.$p$) of $a(\cdot)$ we have, for $j \in \mathbb{Z}^n \setminus 4 \lfloor \sqrt{n} \rfloor Q_{k^0}$, that
\[
(\Phi_{t}^{d} \ast I_{\alpha} \, a)(j) = \sum_{k : |j-k|_{\infty} \leq \frac{t}{4}} \Phi^{d}_{t}(j-k) \sum_{i \in Q_{k^0}} a(i) 
\left[|i-k|^{\alpha-n} - q_{N}(i,k) \right],
\]
where $q_{N}(\, \cdot \,, k)$ is the degree $N - 1$ Taylor polynomial of the function $x \rightarrow |k - x|^{\alpha-n}$ expanded around 
$k^0$. By the standard estimate of the remainder term in the Taylor expansion there exists $\xi$ between $i$ and $k^0$ such that
\[
\left| |i-k|^{\alpha-n} - q_{N}(i,k) \right| \leq C |i - k^0 |^{N} |k - \xi|^{\alpha-n-N}.
\]
Since $|k - k^0 |_{\infty} \geq 2|i - k^0 |_{\infty}$ for any $i \in Q_{k^0}$, we get 
$|k - \xi| \geq |k - \xi|_{\infty} \geq |k - k^0|_{\infty}/2$. So,
\[
\left| |i-k|^{\alpha-n} - q_{N}(i,k) \right| \leq C |i - k^0 |^{N}_{\infty} |k - k^0|^{\alpha-n-N}_{\infty}.
\]
As also $|k - k^0 |_{\infty} \geq |j - k^0 |_{\infty}/2$ for any $j \notin 4\lfloor \sqrt{n} \rfloor Q_{k^0}$, it follows
\[
\left| |i-k|^{\alpha-n} - q_{N}(i,k) \right| \leq C |i - k^0 |^{N}_{\infty} |j - k^0|^{\alpha-n-N}_{\infty},
\]
with $C$ independent of $i$, $j$, $k$ and $k^0$. This inequality, the condition (a2.$p$) of the atom $a(\cdot)$ and the fact that 
$\{ k : |k|_{\infty} \leq t/4 \} \subset \{ k : |k|_{\infty} \leq \lfloor t \rfloor \}$, where $1 \leq \lfloor t \rfloor \leq t$, allow us to conclude that
\[
|(\Phi^{d}_{t} \ast I_{\alpha}a)(j)| \leq C \frac{(2m+1)^{n+N}}{(\# Q_{k^0})^{1/p}} | j - k^0|^{\alpha-n-N}_{\infty} 
\left( \lfloor t \rfloor^{-n} \sum_{|k|_{\infty} \leq \lfloor t \rfloor} |\Phi(k/t)| \right)
\]
\[
\leq C \| \Phi \|_{\infty} \frac{(2m+1)^{n+N}}{(\# Q_{k^0})^{1/p}} | j - k^0|^{\alpha-n-N}_{\infty}
\]
\[
\leq \frac{C}{(\# Q_{k^0})^{1/p}} \left( \frac{(2m+1)^n}{(| j - k^0|_{\infty}^{n})^{1-\frac{\alpha}{n+N}}} \right)^{\frac{n+N}{n}}
\]
\[
\leq \frac{C}{(\# Q_{k^0})^{1/p}} \left( \frac{1}{(\# Q_{j, 2| j - k^0|_{\infty}})^{1-\frac{\alpha}{n+N}}} 
\sum_{i \in Q_{j, 2| j - k^0|_{\infty}}} \chi_{Q_{k^0}}(i) \right)^{\frac{n+N}{n}}
\]
where $Q_{j, 2| j - k^0|_{\infty}}$ is the discrete cube centered at $j$ and side length $4| j - k^0|_{\infty} + 1$. Thus
\[
|(\Phi^{d}_{t} \ast I_{\alpha}a)(j)| \leq \frac{C}{(\# Q_{k^0})^{1/p}} 
\left[ M_{\frac{\alpha n}{n+N}} (\chi_{Q_{k^0}})(j) \right]^{\frac{n+N}{n}}
\]
for all $j \notin 4 \lfloor \sqrt{n} \rfloor Q_{k^0}$ and $1 \leq t \leq 2 |j - k^0 |_{\infty}$. 

\

{\bf Case II:} $|j - k^0 |_{\infty} > 4 \lfloor \sqrt{n} \rfloor m$ and $t > 2 |j - k^0 |_{\infty}$. We have that
\[
(\Phi_t^d \ast I_{\alpha}a)(j) = \sum_{l : |j-l|_{\infty} \leq \frac{t}{4}} \Phi^{d}_{t}(j-l) \sum_{i \in Q_{k^0}} a(i) \, |i-l|^{\alpha-n}
\]
\[
= \sum_{i \in Q_{k^0}} a(i) \sum_{l : |j-l|_{\infty} \leq \frac{t}{4}} \Phi^{d}_{t}(j-l) \, |i-l|^{\alpha-n}
\]
\[
= \sum_{i \in Q_{k^0}} a(i) \sum_{k} \Phi^{d}_{t}(j- k -i) \, |k|^{\alpha-n},
\]
where $k = l - i$, in this case $|j-k-i|_{\infty} \leq t/4$ implies that the inner sum is vanishing on $|k|_{\infty} > 2t$, so
\[
(\Phi_t^d \ast I_{\alpha}a)(j) = \sum_{0 < |k|_{\infty} \leq 2t} |k|^{\alpha-n} \sum_{i \in Q_{k^0}} a(i) \, \Phi^{d}_{t}(j- k -i).
\]
As before, we put $N - 1 = \lfloor n(p^{-1} - 1) \rfloor$. Then, for $j \in \mathbb{Z}^n \setminus 4 \lfloor \sqrt{n} \rfloor Q_{k^0}$, the moment condition (a3.$p$) of $a(\cdot)$ gives
\[
(\Phi_t^d \ast I_{\alpha}a)(j) = \sum_{0< |k|_{\infty} \leq 2t} |k|^{\alpha-n} \sum_{i \in Q_{k^0}} a(i) \, 
\left[\Phi^{d}_{t}(j- k -i) - \widetilde{q}_N(i, j, k, t) \right],
\]
where $\widetilde{q}_N(\, \cdot \,, j, k, t)$ is the degree $N - 1$ Taylor polynomial of the function $x \rightarrow \Phi_{t}(j- k -x)$ expanded around $k^0$. By the standard estimate of the remainder term in the Taylor expansion we have that
\[
\left| \Phi^{d}_{t}(j- k -i) - \widetilde{q}_N(i, j, k, t) \right| \leq C \| \Phi \|_{\mathcal{S}(\mathbb{R}^n), N} \, |i - k^0|^{N} t^{-n-N},
\]
where $C$ does not depend on $i$, $j$, $k$ and $t$. This inequality and the condition (a2.$p$) of the atom $a(\cdot)$ lead to
\[
\left|(\Phi_t^d \ast I_{\alpha}a)(j)\right| \leq C \frac{(2m+1)^{n+N}}{(\# Q_{k^0})^{1/p}}t^{-n-N}\sum_{0 < |k|_{\infty} \leq 2t} |k|^{\alpha-n}.
\]
Now, the estimate (\ref{kernel estimate}) and $t > 2 |j - k^0 |_{\infty}$ imply
\[
\left|(\Phi_t^d \ast I_{\alpha}a)(j)\right| \leq C \frac{(2m+1)^{n+N}}{(\# Q_{k^0})^{1/p}} |j - k^0 |_{\infty}^{\alpha-n-N},
\]
\[
\leq \frac{C}{(\# Q_{k^0})^{1/p}} 
\left[ M_{\frac{\alpha n}{n+N}} (\chi_{Q_{k^0}})(j) \right]^{\frac{n+N}{n}},
\]
for all $j \notin 4 \lfloor \sqrt{n} \rfloor Q_{k^0}$ and $t > 2 |j - k^0 |_{\infty}$. 

Thus, (\ref{cero}) and the estimates obtained in the cases $I$ and $II$ above give
\begin{equation} \label{cota afuera}
\sum_{j \in \mathbb{Z}^n \setminus 4\lfloor \sqrt{n} \rfloor Q_{k^0}} \sup_{t>0}|(\Phi^{d}_{t} \ast I_{\alpha}a)(j)|^q \leq 
\frac{C}{(\# Q_{k^0})^{q/p}} \sum_{j \in \mathbb{Z}^n} \left[ M_{\frac{\alpha n}{n+N}} (\chi_{Q_{k^0}})(j) \right]^{q \frac{n+N}{n}}.
\end{equation}
Since  $N-1= \lfloor n(\frac{1}{p}-1) \rfloor$, we have $q \frac{n+N}{n} > 1$. We write $\widetilde{q} = q \frac{n+N}{n}$ and let 
$\frac{1}{\widetilde{p}} = \frac{1}{\widetilde{q}} + \frac{\alpha}{n+N}$, so $\frac{\widetilde{p}}{\widetilde{q}} = \frac{p}{q}$. From
Proposition 2.4 in \cite{PRocha}, we obtain
\begin{equation} \label{cota afuera 2}
\sum_{j \in \mathbb{Z}^n} \left[ M_{\frac{\alpha n}{n+N}} (\chi_{Q_{k^0}})(j) \right]^{q \frac{n+N}{n}} \leq
C \left( \sum_{j \in \mathbb{Z}^n} \chi_{Q_{k^0}}(j) \right)^{q/p} = C (\# Q_{k^0})^{q/p}.
\end{equation}
Now, (\ref{cota afuera}) and (\ref{cota afuera 2}) give
\begin{equation} \label{cota afuera 3}
\sum_{j \in \mathbb{Z}^n \setminus 4\lfloor \sqrt{n} \rfloor Q_{k^0}} \sup_{t>0}|(\Phi^{d}_{t} \ast I_{\alpha}a)(j)|^q \leq C.
\end{equation}
Finally, (\ref{estim C}) and (\ref{cota afuera 3}) lead to (\ref{uniform estimate}). Thus the proof is finished.
\end{proof}

\begin{theorem} \label{main result}
For $0 < \alpha < n$, let $I_{\alpha}$ be the discrete Riesz potential given by (\ref{Riesz potential}). If
$0 < p \leq 1$ and $\frac{1}{q} = \frac{1}{p} - \frac{\alpha}{n}$, then
\[
\| I_{\alpha} \, b \|_{H^{q}(\mathbb{Z}^n)} \leq C \| b \|_{H^{p}(\mathbb{Z}^n)}
\]
for all $b \in H^{p}(\mathbb{Z}^n)$, where $C$ does not depend on $b$.
\end{theorem}

\begin{proof}
We fix $p_0$ such that $1 < p_0 < \frac{n}{\alpha}$ and $\Phi \in \mathcal{S}(\mathbb{R}^n)$ with $\supp(\Phi) \subset B_{1/4}({\bf 0 })$ 
and $\int \Phi = 1$. By Theorem \ref{atomic Hp}, given $b \in H^{p}(\mathbb{Z}^n)$, with $0 < p \leq 1$, we can write 
$b = \sum_k \lambda_k a_k$ where the $a_k$'s are $(p, \infty, N_p)$ atoms, the scalars $\lambda_k$ satisfies 
$\sum_{k} |\lambda_k |^{p} \leq C \| b \|_{H^{p}(\mathbb{Z}^n)}^{p}$  and the series converges in $H^{p}(\mathbb{Z}^n)$ and 
so in $\ell^{p_0}(\mathbb{Z}^n)$ since $H^{p}(\mathbb{Z}^n) \subset \ell^{p}(\mathbb{Z}^n) \subset \ell^{p_0}(\mathbb{Z}^n)$ embed continuously. Now for $\frac{1}{q_0} = \frac{1}{p_0} - \frac{\alpha}{n}$, by \cite[Theorem 3.1]{PRocha}, the discrete Riesz potential 
$I_{\alpha}$ is a bounded operator $\ell^{p_0}(\mathbb{Z}^n) \to \ell^{q_0}(\mathbb{Z}^n)$ and since 
$H^{q_0}(\mathbb{Z}^{n}) = \ell^{q_0}(\mathbb{Z}^{n})$ with equivalent norms (see Remark \ref{Hp equiv}), it follows that
\begin{equation} \label{discrete maximal}
\sup_{t>0} |(\Phi_{t}^d \ast I_{\alpha}b)(j)| \leq \sum_{k=1}^{\infty} |\lambda_k| \sup_{t>0}|(\Phi_{t}^d \ast I_{\alpha} a_k)(j)|, \,\,\, \text{for all} \,\, j \in \mathbb{Z}^n.
\end{equation}
Then for $1 \leq q \leq \frac{n}{n - \alpha}$, by (\ref{discrete maximal}) and Minkowski's integral inequality on $\sigma$-finite measure spaces, we have
\begin{equation} \label{mink ineq}
\|\sup_{t>0} |(\Phi_{t}^d \ast I_{\alpha}b)| \|_{\ell^{q}(\mathbb{Z}^n)} \leq \sum_{k} |\lambda_k| \,
\| \sup_{t>0}|(\Phi_{t}^d \ast I_{\alpha} a_k)| \|_{\ell^{q}(\mathbb{Z}^n)}.
\end{equation}
Now for $0 < q < 1$, from (\ref{discrete maximal}), it is easy to check that
\begin{equation} \label{q ineq}
\| \sup_{t>0} |(\Phi_{t}^d \ast I_{\alpha}b)| \|_{\ell^{q}(\mathbb{Z}^n)}^q \leq \sum_{k} |\lambda_k|^q 
\| \sup_{t>0}|(\Phi_{t}^d \ast I_{\alpha} a_k)| \|_{\ell^{q}(\mathbb{Z}^n)}^q.
\end{equation}
Finally, since $0 < p \leq 1$ and $0 < p < q \leq \frac{n}{n - \alpha}$, the estimate (\ref{mink ineq}) or (\ref{q ineq}) according to the case, Proposition \ref{maximal atom estim}, and the fact that $\sum_{k} |\lambda_k |^{p} \leq C \| b \|_{H^{p}(\mathbb{Z}^n)}^{p}$  allow us to conclude that 
\[
\|\sup_{t>0}|(\Phi_{t}^d \ast I_{\alpha} b)| \|_{\ell^{q}(\mathbb{Z}^n)} \leq C \left( \sum_{k} |\lambda_k|^{\min\{1, q \}} \right)^{\frac{1}{\min\{1, q \}}} \leq C \left( \sum_{k} |\lambda_k |^{p} \right)^{1/p} \leq C\| b \|_{H^{p}(\mathbb{Z}^n)}.
\] 
Since $b$ is an arbitrary element of $H^{p}(\mathbb{Z}^n)$, the theorem follows.
\end{proof}

\begin{remark}
The argument used to prove Theorem 3.2 is similar to the one given in the proof of \cite[Theorem 2.4.5, see p. 140-141]{Grafakos}.
\end{remark}

{\bf Acknowledgements.} I am very grateful to the referee for the careful reading of the paper and for the useful comments and suggestions which helped me to improve the original manuscript.

\begin{thebibliography}{99}

\bibitem{Bownik} M. Bownik, {\it Boundedness of operators on Hardy spaces via atomic decompositions}, Proc. Amer. Math. Soc., 133 (2005), 
3535-3542.

\bibitem{Boza} S. Boza and M. Carro, {\it Discrete Hardy spaces}, Studia Math., 129 (1) (1998), 31-50.

\bibitem{Carro} S. Boza and M. Carro, {\it Hardy spaces on $\mathbb{Z}^N$}, Proc. R. Soc. Edinb., 132 A (1) (2002), 25-43.

\bibitem{Coif} R. Coifman, {\it A real variable characterization of $H^{p}$}, Studia Math., 51 (1974), 269-274.

\bibitem{Feff} C. Fefferman and E. Stein, {\it $H^p$ spaces of several variables}, Acta Math. 129 (1972), 137-193.

\bibitem{Loukas} L. Grafakos, Classical Fourier Analysis, 3rd ed., Graduate Texts in Mathematics 249, Springer New York, 2014.

\bibitem{Grafakos} L. Grafakos, Modern Fourier Analysis, 3rd ed., Graduate Texts in Mathematics 250, Springer New York, 2014.

\bibitem{Hardy} G. H. Hardy, J. E. Littlewood and G. P\'olya, Inequalities, 2nd ed., Cambridge Univ. Press, London and new York, 1952.

\bibitem{Kanjin} Y. Kanjin and M. Satake, {\it Inequalities for discrete Hardy spaces}, Acta Math. Hungar., 89 (4) (2000), 301-313.

\bibitem{Krantz} S. Krantz, {\it Fractional integration on Hardy spaces}, Studia Mathematica, vol 73 (2) (1982), 87-94.

\bibitem{Latter} R. Latter, {\it A characterization of $H^{p}(\mathbb{R}^{n})$ in terms of atoms}, Studia Math., 62 (1978), 93-101.

\bibitem{Nakai} E. Nakai, {\it Recent topics of fractional integrals}, Sugaku Expo. 20, No. 2 (2007), 215-235. 

\bibitem{Ricci} F. Ricci and J. Verdera, {\it Duality in spaces of finite linear combinations of atoms}, Trans. Amer. Math. Soc. 363 
(2011), 1311-1323.

\bibitem{Pablo} P. Rocha, {\it A note on Hardy spaces and bounded linear operators}, Georgian Math. J., 25 (1) (2018), 73-76.

\bibitem{Rocha} P. Rocha, {\it Fractional series operators on discrete Hardy spaces}, Acta Math. Hungar., 168 (1) (2022), 202-216.

\bibitem{PRocha} P. Rocha, {\it A note about discrete Riesz potential on $\mathbb{Z}^n$}. To appear in Publ. Math. Debr.

\bibitem{Rocha2} P. Rocha, {\it A molecular decomposition for $H^p(\mathbb{Z}^n)$}, J. Class. Anal., vol 26 (1) (2025), 63-76.

\bibitem{Stein} E. Stein, Singular Integrals and Differentiability Properties of Functions, Princeton Univ. Press, 1970.

\bibitem{St} E. Stein, Harmonic Analysis: Real-Variable Methods, Orthogonality, and Oscillatory Integrals, Princeton Univ. Press,
1993.

\bibitem{Wainger} E. Stein and S. Wainger, {\it Discrete analogues in harmonic analysis. II: Fractional integration}, J. Anal. Math. 80 
(2000), 335-355.

\bibitem{Weiss} E. Stein and G. Weiss, {\it On the theory of harmonic functions of several variables}, Acta Math. 103 (1960), 25-62.

\bibitem{Taible} M. H. Taibleson and G. Weiss, {\it The molecular characterization of certain Hardy spaces}, Asterisque 77 (1980), 67-149.


\end{thebibliography}

Pablo Rocha, Instituto de Matem\'atica (INMABB), Departamento de Matem\'atica, Universidad Nacional del Sur (UNS)-CONICET, Bah\'ia Blanca, Argentina. \\
{\it e-mail:} pablo.rocha@uns.edu.ar

\end{document} 
